\documentclass{article}%Do not change this

\usepackage{graphicx}

\title{Cyberpunk 2077: Ethics of Connected Technology}
\author{Quinton Baudoin, Caylei Dauphiney, Royce Diano, Connor Segura}

\date{Fall 2026} %Leave as-is.

\begin{document}
\maketitle

\begin{abstract}
\textit{Cyberpunk 2077} depicts a society where corporations deploy powerful networked systems while shifting their risks onto the public. This report examines three connected issues: accountability for autonomous artificial intelligence (AI), safety and equitable access to cybernetic medical devices, and hacking that threatens privacy and physical infrastructure. For each, we identify the harm, affected parties, the story's response, and practical safeguards. The game's exaggerated technology highlights a current lesson for computing professionals: security, oversight, and access must be built into systems rather than left to individual users or addressed only after harm occurs.
\end{abstract}

\section{Introduction}
\textit{Cyberpunk 2077}, developed by CD Projekt Red, takes place in Night City, where corporations wield influence that often exceeds public institutions. The player character, V, becomes involved in conflicts over corporate technology and digital identity. Because the plot is readily available, this report gives only the context needed to analyze its ethical issues. We focus on how the game portrays systems that affect people's safety, bodies, privacy, and access to essential services.

\section{Artificial Intelligence and Accountability}
\subsection{The Datakrash and rogue AI}
Before the Datakrash, corporations and governments developed powerful AI systems, including experimental programs and military-grade weapons, and kept some confined within controlled parts of the Old Net. The netrunner Rache Bartmoss released destructive programs known as R.A.B.I.D.S. into that network. The resulting Datakrash corrupted data and disrupted the safeguards containing dangerous AIs, leaving rogue intelligences free to operate in the ruined Old Net. They are not all identical: some were designed for hostile purposes, while others may have changed or developed beyond their creators' control.

The ethical problem begins before the escape. Organizations built potentially dangerous systems in secret, but the consequences of losing control were borne by society, not just their creators. Bartmoss's attack then compounded that danger by damaging shared infrastructure and exposing people to threats they had not chosen. Netrunners, users of connected services, and communities dependent on communications all face risks, while the original developers and attackers may be difficult to hold accountable.

The Datakrash illustrates the limits of relying on secrecy or a sandbox alone: confinement can fail if its surrounding infrastructure is attacked. Before deploying powerful systems, developers should independently assess foreseeable misuse, restrict dangerous capabilities, test containment and recovery, and establish responsibility for harm. These safeguards cannot eliminate every risk, but they make safety a shared obligation rather than a problem deferred until after a catastrophe. The Blackwall is the setting's later attempt to manage what escaped.

\subsection{The Blackwall and limits of containment}
After the Datakrash, the Blackwall separates the surviving network from rogue AIs in the Old Net. It is not just background lore: V's search for a way to remove the Relic leads the Voodoo Boys to use V to reach Alt Cunningham, an intelligence beyond the barrier. Their plan turns a life-threatening breach into a tool for their community's survival, while NetWatch tries to stop access and protect the wider network. V is caught between groups with competing security and political goals and is exposed to risks that neither group meaningfully lets V evaluate. The episode makes the Blackwall central to the main story's conflict over control of digital identity and the power to cross boundaries between human and machine.

The Blackwall is not the setting's only defense: city networks are isolated from one another, with limited, controlled communication between them. Nevertheless, powerful actors continue to probe the barrier for strategic advantage. A breach could endanger network users and communities who did not choose to accept that risk, while decisions are made by secretive corporate and security institutions. The issue is also ethical, not only technical. Calling every intelligence beyond the wall a hostile "rogue AI" may obscure differences between a destructive program, a system acting in self-defense, and a digital person such as Alt, whose identity originated as a human engram. The story leaves unresolved who has the authority to decide whether such entities have rights or may be destroyed.

NetWatch presents itself as a guardian of the barrier, but its power and operations receive little public accountability or independent review. The Voodoo Boys likewise pursue their aims without giving V meaningful information or consent. Existing network isolation reduces exposure, but it does not resolve who controls access, who may take risks with the Old Net, or who is accountable when others are harmed. A better approach would build on those protections with independent oversight, clear limits on high-risk access, tested recovery plans, and informed consent for people exposed to experiments. Affected communities should have a voice in decisions, rather than leaving them solely to security agencies and powerful organizations.

\section{Cyberware, Health, and Bodily Safety}
\subsection{Corporate control and medical access}
Cyberware ranges from elective enhancements to devices that support bodily functions. The Corpo prologue makes dependence personal: after Arasaka fires V, it disables V's implants along with corporate access. Losing a job also means losing control over technology integrated with one's body, giving an employer leverage over a former worker.

Access to upgrades and care in Night City depends on wealth; people unable to afford reputable service may turn to unsafe clinics. Companies can also control repairs, replacements, software support, or even disconnection. Patients and lower-income residents are most exposed, yet the game treats this imbalance as normal. A fairer system would guarantee support and repair for medically necessary devices and prevent employers or vendors from disabling essential bodily functions. Interoperability and right-to-repair protections would limit dependence on a single company.

\subsection{Cybersecurity and weaponization}
The game depicts quickhacks disabling or taking control of connected implants. Although dramatized, the underlying concern is credible: when software controls a medical device, a security failure can become a safety failure. Patients are directly affected, and nearby people may also be put at risk if a device is misused. The setting mostly leaves protection to individuals who can afford better defenses.

The U.S. Food and Drug Administration recommends that medical-device makers address cybersecurity throughout a product's life cycle, including secure design and management of emerging risks \cite{FDA2023}. NIST's IoT device baseline similarly identifies core capabilities such as protecting device data, controlling interfaces, and supporting software updates \cite{NIST2020}. These measures should be minimum manufacturer responsibilities, not premium features: devices need secure update paths, limited privileges, vulnerability reporting, and support plans that account for long-lived implants.

The story also shows medical enhancements repurposed as weapons. This dual use can harm bystanders and users; cyberpsychosis raises concerns about psychological consequences and access to care. Augmentation should not be treated as an inevitable cause of violence. Night City's response is largely force after a crisis. Prevention would include informed consent, safe installation, mental-health care, and proportionate limits on weapon-grade modifications, without assuming that disability makes a person dangerous.

\section{Hacking, Privacy, and Infrastructure}
\subsection{Surveillance and unauthorized access}
The Ping quickhack reveals nearby devices and people, enabling surveillance without their knowledge or consent. It can expose private behavior and enable targeting, including of people who never chose to connect. Collected information can also be reused to track or discriminate. In the story, privacy is left to individuals while corporations and criminals exploit weak security; people being watched may not know or be able to challenge its use.

Real devices also need safeguards against unauthorized discovery and access. NIST's baseline emphasizes limiting access to device interfaces and protecting data \cite{NIST2020}. Systems should minimize collected information, use secure defaults, and provide clear notice and meaningful control. Organizations should limit who can query personal data, retain it only as long as needed, and be accountable for misuse. Security alone does not make surveillance ethical: collection must also be necessary, proportionate, and governed by consent and law.

\subsection{Physical harm and resilient services}
Night City's netrunners use digital access to cause bodily injury and disrupt services. Such attacks cross from data loss into physical harm; patients, workers, bystanders, and communities relying on essential services may all be affected. The Datakrash shows how disruption to communication systems can cascade across a digitally dependent society. Yet hacking is treated as routine conflict, with little accountability for attacks on civilians.

The response should combine prevention, resilience, and accountability. Organizations should separate critical systems, restrict access, monitor incidents, and rehearse recovery using safe fallback procedures. Offline or manual modes can prevent outages from becoming public-safety crises. Makers and operators should coordinate vulnerability handling and disclose incidents. These measures reduce cascading harm; responsibility for foreseeable risks should not fall solely on victims.

\section{Conclusion}
Across AI, cyberware, and hacking, \textit{Cyberpunk 2077} shows technology deployed without adequate oversight, leaving people to bear risks they did not choose. The details are fictional, but the obligations are practical: build systems to fail safely, protect users throughout a product's life, ensure equitable access and repair, and remain answerable to affected people.

\bibliography{citations}
\bibliographystyle{ieeetr}

\end{document}
